\section{结论}

本报告基于 Li 与 Hu（2014）的论文框架，使用 MPI 并行编程模型完整实现了三种层次递增的分布式进化算法——Ring DEA、Ring Individual HDEA 和 Ring Colony HDEA（HdEA-MP），并以 TSPLIB 实例 pcb442 为测试平台，在 64 子种群（8×8 组）、严格遵循论文参数配置（Inver-over 算子、映射算子、临界速度机制、变异率衰减、Ring 拓扑、20:1 局部-全局比、100 全局迁移轮次）的条件下进行了系统的 10 次对比实验。

实验结果表明，两类层次化算法均以极高的统计显著性（$p<0.001$）优于基础分布式算法 dEA，其中移动种群策略的 HdEA-MP 在均值（50972）、标准差（76）和单次最优值（50821）三项指标上均为三者最优，且其运行速度约为 dEA 的两倍（347 秒对 719 秒）。层级提升的求解质量呈现出清晰的递增梯度：dEA（51409）→ HdEA（51053）→ HdEA-MP（50972），每一步升级均为求解质量带来了可量化的改善。HdEA-MP 相对于 HdEA 的优势在当前 10 次运行的样本量下尚未达到 0.05 的统计显著水平（$p=0.115$），但其在各项指标上的全面领先预示着在更大的计算规模下这一优势极有可能被统计检验所证实。

从架构层面来看，层次化迁移结构的核心价值在于：通过"微观频繁、宏观稀疏"的双层节奏设计，在子种群之间的基因交流与遗传隔离之间找到了一个平衡点，使得并行架构的独立探索优势与群体智能的信息共享需求能够同时被满足，而不会像 dEA 的扁平拓扑那样因过度交流而导致多样性过早丧失。移动种群策略则进一步深化了这一设计——将单次全局迁移的信息注入从"一个体"的量级跃升至"一个完整子种群"的量级，以边际的成本换取了质的飞跃。

本实验受限于单机 20 核的计算资源，未能以论文中的计算量（200M 代，1264 核集群）进行完整复现——两者的总计算量相差约 333 倍，这导致了本实验的绝对求解精度尚不及论文。然而，算法排名的一致性和统计显著性的稳健性充分证明了：论文中提出的架构的核心价值并不依赖于超大规模的计算资源——即使在有限资源条件下，算法设计的优劣依然是决定求解质量的主导因素。在未来的工作中，可以在更多 TSPLIB 实例（尤其是论文中的高难度实例如 u724、dsj1000 等）上扩展实验，并将子种群数扩展至论文的 64 独立物理核心配置，以获得与原始论文完全可比的求解精度。
